% Propietario conservado: /Users/ruben/Documents/ChatGPT/jueces y controles/REVISION_ARGUMENTAL_APERTURAS_CIERRES_20260915/delivery/ARTICLE_V_ES_ARGUMENTAL_20260915/manuscrito/sections/acoplamiento_recuperacion.tex
\section{Couplage arithmético-géométrique et récupération bilatérale}
\label{v:sec:acoplamiento-recuperacion}

Le registre dodécaphasique de la section~\ref{v:sec:k-registro} admet
des lectures arithmétiques, d’incidence et opératorielles. Nous appelons
\emph{constante holographique de couplage arithmético-géométrique}
le registre engendré \(K\), en conservant l’ordre de ses douze fenêtres,
l’origine de phase et l’orientation. Sa lecture rationnelle
\(\kappa=K_{\mathrm{int}}/(1000^{12}-1)\) est un codage de ce
registre ; elle ne remplace ni ses types ni l’histoire enrichie qui le produit.
L’inversion de Hadamard récupère sa coordonnée signée ; la lecture
d’incidence détermine la tétrade de Witt. Les constructions suivantes
explicitent deux autres fonctions : identifier des opérateurs sur son porteur
cyclique et conserver de l’information entre deux réalisations isométriques.
Ce sont des opérations postérieures à la réalisation du continu, non des règles
qui introduisent des constantes conventionnelles dans APP--TRIT--TPK.

\subsection{Codage exact et référence cyclique}
\label{v:subsec:acoplamiento-codigo}

\begin{proposition}[Récupération du registre à partir de sa lecture rationnelle]
La base \(b=1000\), la longueur douze, l’origine et l’orientation étant fixées, l’application
\[
 (k_1,\ldots,k_{12})\longmapsto
 \kappa=\frac{\sum_{j=1}^{12}k_jb^{12-j}}{b^{12}-1},
 \qquad 0\leq k_j<b,
\]
est injective. À partir de \(\kappa\), on récupère l’entier
\(n=(b^{12}-1)\kappa\) et, par divisions euclidiennes successives, les douze
blocs, y compris les zéros initiaux. Une erreur absolue strictement
inférieure à \(1/[2(b^{12}-1)]\) permet la même récupération par arrondi.
\end{proposition}
\begin{proof}
Le développement entier de longueur fixe en base \(b\) est unique. Deux entiers
distincts sont séparés d’au moins un ; leurs lectures rationnelles
sont séparées d’au moins \(1/(b^{12}-1)\). La borne stricte situe la
donnée approchée à l’intérieur de la cellule d’arrondi d’un unique entier.
\end{proof}

Considérons maintenant \(K=(234,543,140,729,659,824,621,58,914,794,146,601)\),
le décalage cyclique \(S\) de ses positions et
\(O_K=(K,SK,\ldots,S^{11}K)\) sur \(\mathbb R^{12}\).
Le décalage des positions n’est identifié ni à toute l’évolution
enrichie ni, par sa seule périodicité, à une symétrie exceptionnelle.

\begin{theorem}[Inversibilité de la référence orbitale]
\label{v:thm:acoplamiento-orbita}
La matrice \(O_K\) est inversible et
\[
 589609\|c\|^2\leq\|O_Kc\|^2\leq6263^2\|c\|^2.
\]
Son orbite centrée est de rang onze.
\end{theorem}
\begin{proof}
La matrice de Gram \(O_K^*O_K\) est circulante. Sa première ligne, obtenue au moyen
des douze produits scalaires \(\langle K,S^jK\rangle\), est
\[
 \begin{aligned}
 (&4251257,2999323,3163385,3287920,3216553,3344743,\\
  &2950064,3344743,3216553,3287920,3163385,2999323).
 \end{aligned}
\]
L’évaluation sur les douze caractères du cycle fournit les valeurs
propres
\[
 \begin{gathered}
 6263^2,\quad697225,\quad
 589609\ (2),\quad1407529\ (2),\quad1053157\ (2),\\
 (1248025-345420\sqrt3)\ (2),\qquad
 (1248025+345420\sqrt3)\ (2),
 \end{gathered}
\]
où les parenthèses indiquent les multiplicités. La plus petite valeur est
\(589609\) : pour la seule comparaison non immédiate,
\(658416>345420\sqrt3\), comme le montre le carré des deux membres.
L’inégalité triangulaire de la transformée de Fourier et \(K_j\geq0\)
donnent le maximum \((\sum K_j)^2=6263^2\).
La positivité de tous les modes prouve l’inversibilité et les bornes.
Centrer le registre élimine uniquement le mode uniforme.
\end{proof}

\begin{corollary}[Identification d’un opérateur par ses réponses]
\label{v:cor:acoplamiento-identificacion}
Un opérateur linéaire \(T:\mathbb R^{12}\to\mathbb R^{12}\) est déterminé par
\(Y=(TK,TSK,\ldots,TS^{11}K)\), au moyen de \(T=YO_K^{-1}\).
Si \(TS=ST\), la réponse vectorielle complète \(y=TK\) suffit, car
\(T=O_yO_K^{-1}\). Pour une perturbation \(E\) de \(Y\),
\(\|\widehat T-T\|\leq\|E\|/\sqrt{589609}\).
\end{corollary}
\begin{proof}
On a \(Y=TO_K\). Sous commutation, \(TS^jK=S^jy\), donc \(Y=O_y\).
La borne se déduit de
\(\|O_K^{-1}\|=1/\sqrt{589609}\).
Les normes sont la norme euclidienne et sa norme d’opérateur subordonnée.
\end{proof}

La réponse \(y\) contient douze composantes. L’affirmation n’attribue pas
à un unique nombre scalaire la récupération d’un opérateur arbitraire :
elle précise quelle observation et quelle commutation rendent l’identification possible.

\subsection{Conservation et restitution de deux canaux}
\label{v:subsec:acoplamiento-bilateral}

Soient \(H,H'\) des espaces de Hilbert et
\(\mathcal B,\mathcal C:H\to H'\) des isométries linéaires bornées.
Pour \(a>0\), nous définissons
\[
 p=a^2+a+1,\qquad N=(2a+1)p=(a+1)^3+a^3,\qquad
 Q_-=a\mathcal B-\mathcal C,\quad Q_+=(a+1)\mathcal B+\mathcal C.
\]

\begin{theorem}[Bilan bilatéral et récupération]
\label{v:thm:acoplamiento-balance}
On a
\[
 (a+1)Q_-^*Q_-+aQ_+^*Q_+=NI_H,\qquad
 W_av=\frac1{\sqrt N}
 \begin{pmatrix}\sqrt{a+1}\,Q_-v\\\sqrt a\,Q_+v\end{pmatrix},
 \qquad W_a^*W_a=I_H.
\]
Les données \(x=Q_-v\), \(y=Q_+v\) déterminent
\[
 \mathcal Bv=\frac{x+y}{2a+1},\qquad
 \mathcal Cv=\frac{ay-(a+1)x}{2a+1},\qquad
 v=\mathcal B^*\frac{x+y}{2a+1}.
\]
\end{theorem}
\begin{proof}
En posant \(X=\mathcal B^*\mathcal C+\mathcal C^*\mathcal B\), on obtient
\[
 Q_-^*Q_-=(a^2+1)I-aX,\qquad
 Q_+^*Q_+=((a+1)^2+1)I+(a+1)X.
\]
Les pondérations annulent \(X\) et le coefficient restant est
\(2a^3+3a^2+3a+1=N\). La normalisation donne l’isométrie.
Additionner les deux canaux récupère \((2a+1)\mathcal Bv\) ;
la combinaison \(ay-(a+1)x\) récupère \((2a+1)\mathcal Cv\).
Enfin, \(\mathcal B^*\mathcal B=I_H\).
\end{proof}

Le bilan vaut aussi en dimension infinie : on ne compose que des
opérateurs bornés. L’image de \(W_a\) est fermée et son projecteur
orthogonal est \(W_aW_a^*\). C’est pourquoi une paire normalisée admet une
entrée exactement lorsqu’elle appartient à cette image ; sa composante
orthogonale mesure le défaut de compatibilité.

La factorisation
\[
 W_a=R_aE_a,\qquad
 E_av=\frac1{\sqrt p}
 \begin{pmatrix}\sqrt{a(a+1)}\,\mathcal Bv\\\mathcal Cv\end{pmatrix},
 \quad
 R_a=\frac1{\sqrt{2a+1}}
 \begin{pmatrix}\sqrt a&-\sqrt{a+1}\\
 \sqrt{a+1}&\sqrt a\end{pmatrix}
\]
se vérifie en multipliant les blocs ; \(R_a\) agit sur
\(H'\oplus H'\), chaque entrée scalaire multipliant \(I_{H'}\).
Dans la partition nonadique
\(a=8\), les poids sont \(72/73\) et \(1/73\) ; le facteur
\(73=8\cdot9+1\) est une conséquence de la loi générale
\(a(a+1)+1\), non une donnée primitive ajoutée. Le bilan correspondant est
\(9\|Q_-v\|^2+8\|Q_+v\|^2=17\cdot73\,\|v\|^2\).

\subsection{Transport des observables et mémoire entrante}
\label{v:subsec:acoplamiento-lectores}

\begin{theorem}[Observables communs aux deux canaux]
\label{v:thm:acoplamiento-lector}
Pour un opérateur borné autoadjoint \(F\) de \(H'\),
\[
 W_a^*\operatorname{diag}(F,F)W_a
 =\frac{a(a+1)\mathcal B^*F\mathcal B+\mathcal C^*F\mathcal C}{p}.
\]
Si \(\mathcal B\) est unitaire et \(\mathcal C=\mathcal BR\), avec
\(R\) unitaire, et \(G=\mathcal B^*F\mathcal B\), l’observable est
conservé si et seulement si \(GR=RG\).
\end{theorem}
\begin{proof}
En développant \((a+1)Q_-^*FQ_-+aQ_+^*FQ_+\), les termes croisés
s’annulent. Les deux autres coefficients sont
\(a(a+1)(2a+1)\) et \(2a+1\). Diviser par \(N\) donne la formule.
Dans le cas indiqué, l’égalité avec \(G\) équivaut à \(R^*GR=G\) ;
l’unitarité la rend équivalente à \(GR=RG\).
\end{proof}

La norme s’obtient avec \(F=I\). Un observable général requiert
la condition de commutation : conserver la norme n’équivaut pas à conserver
chaque observable sans transformation.

\begin{theorem}[Extension unitaire avec mémoire]
\label{v:thm:acoplamiento-memoria}
Si \(H'=H\), \(\mathcal B=I\) et que \(\mathcal C=U\) est unitaire, soient
\[
 A=\frac{\sqrt{a+1}}{\sqrt N}Q_-,\qquad
 D=\frac{\sqrt a}{\sqrt N}Q_+,\qquad
 \mathscr R_a=\begin{pmatrix}A&-D^*\\D&A^*\end{pmatrix}.
\]
Alors \(\mathscr R_a\) est unitaire dans \(H\oplus H\), sa première
colonne est \(W_a\), et il conserve \(\|v\|^2+\|m\|^2\) pour toute
entrée \((v,m)\). L’inversion est
\[
 v=A^*z_1+D^*z_2,\qquad m=-Dz_1+Az_2.
\]
\end{theorem}
\begin{proof}
Les opérateurs \(A,D\) sont normaux et commutent entre eux et avec leurs
adjoints, étant des polynômes de \(U,U^*\). Le bilan donne
\(A^*A+D^*D=I\), et la normalité donne \(AA^*+DD^*=I\).
Les termes croisés des deux produits
\(\mathscr R_a^*\mathscr R_a\) et \(\mathscr R_a\mathscr R_a^*\)
s’annulent ; les termes diagonaux sont l’identité. Appliquer l’adjoint à
\((z_1,z_2)\) fournit l’inverse.
\end{proof}

La récupération du registre, l’identification orbitale et l’extension
unitaire précisent trois niveaux d’information distincts. Le premier
conserve les blocs avec leur carte ; le deuxième identifie des opérateurs sous
les hypothèses déclarées ; le troisième transporte conjointement l’état et la
mémoire. Cette hiérarchie permet d’utiliser la même structure générée sans
confondre une lecture réduite avec la totalité de l’état.
